% This is a Template for your IPA Documentation

\documentclass{article}
\usepackage[utf8]{inputenc}
\usepackage[german]{babel}

\usepackage{multirow}
\usepackage{tabularx}
\usepackage{subfig}
\usepackage{graphicx}

\usepackage{cite}

\usepackage{todonotes}

\usepackage[utf8]{inputenc}
\usepackage[acronym]{glossaries}

% Add dotted lines to TOC
\usepackage{tocloft}

%---------Tree:Begin-----------------
\usepackage{tikz}
%%%<
\usepackage{verbatim}
%%%>
\usetikzlibrary{trees}
%----------Tree:End------------------

%----------Glossary:Start------------
\makeglossaries

% Sample glossary entry
\newglossaryentry{maths}
{
    name=mathematics,
    description={Mathematics is what mathematicians do}
}

% Sample acronym
\newacronym{mt}{MT}{Math Teachers}

%----------Glossary:End--------------

%----------Title:Start---------------
\title{The \LaTeX\  IPA Template}
\author{Max Muster}
\date{1. Mai 2020}
%\institute{Quality \& Quantity GmbH, Musterstrasse 12, 0000 Stadt}
%----------Title:End-----------------

\begin{document}

\maketitle

\newpage
\begin{table}
    \begin{tabularx}{0.5\textwidth}{l l}
        \textbf{Autor} & Max Muster \\
        \hline
        \textbf{Datum} & 1. Mai 2020 \\
        \hline
        \textbf{Thema} & IPA Titel \\
        \hline
    \end{tabularx}
\end{table}
\begin{table}
    \begin{tabularx}{1.0\textwidth}{l l l}
        \multirow{3}{*}{\textbf{Geht an}} & Hauptexperte & Name Expert 1 \\
        \cline{2-3}
        & Nebenexperte & Name Expert 2 \\
        \cline{2-3}
        & Verantwortliche Fachkraft & Name VeFa \\
        \cline{1-3}
        \textbf{Zur Kenntnis} & Berufsbildner & Name Berufsbildner \\
        \cline{1-3}
    \end{tabularx}
\end{table}
\begin{table}
    \begin{tabularx}{1.0\textwidth}{l l}
        \textbf{Version} & 10.0 \\
        \hline
        \textbf{Ersetzt Version} & - \\
        \hline
        \textbf{Ausgabedatum} & 1. Mai 2020 \\
        \hline
        \textbf{Gültig ab} & 1. Mai 2020\\
        \hline
        \textbf{Gültig bis} & - \\
        \hline
        \textbf{Dokumentname} & \jobname.tex \\
        \hline
    \end{tabularx}
\end{table}

\clearpage

\tableofcontents

%Start of Report
\newpage
\part{Umfeld und Ablauf}
\newpage
\section{Zielgruppe}
\Gls{maths} is very Important for \acrfull{mt}
\section{Aufgabenstellung (gemäss PkOrg)}
\begin{itemize}
    \item{Wichtige Anforderung 1}
    \item{Wichtige Anforderung 1}
    \item{Wichtige Anforderung 1}
    \item{Nicht-Funktionale-Anforderung 1}
\end{itemize}
\section{Deklaration der Firmenstandards}
\section{Projektaufbauorganigramm}
\section{Organisation}
\section{Ablauf}
\section{Meilensteine}
\section{Zeitplan}
\section{Arbeitsjournal}
\begin{table}[h]
    \centering
    \begin{tabular}{l | l | l}
    A & B & C \\
    \hline
    1 & 2 & 3 \\
    4 & 5 & 6
    \end{tabular}
    \caption{very basic table}
    \label{tab:abc}
    \end{table}
\section{Projektjournal}

\newpage
\part{Projekt-Dokumentation}
\newpage
\section{Kurzfassung}
\section{Initialisierung}
In Abb. \ref{fig:me_working} sieht man mich arbeiten. Abb. \ref{fig:me_working:b} zeigt meinen Arbeitsplatz sehr schön.
\begin{figure}[h!]
    \subfloat[Ich beim Arbeiten]{\label{fig:me_working:a}\includegraphics[width=0.3\columnwidth]{figures/hard_working_1.jpg}}
    \hfill
    \subfloat[Ich beim Arbeiten]{\label{fig:me_working:b}\includegraphics[width=0.3\columnwidth]{figures/hard_working_2.jpg}}
    \hfill
    \subfloat[Ich beim Arbeiten]{\label{fig:me_working:c}\includegraphics[width=0.3\columnwidth]{figures/hard_working_3.jpg}}
    \caption{Drei Beispiele wie ich gearbeitet habe.}
    \label{fig:me_working}
\end{figure}
\section{Konzeptphase}
\section{Realisierung}
\begin{figure}[h!]
    \centering
    \includegraphics[width=0.4\columnwidth]{figures/my_product.jpg}
    \caption{Ein Bild meiner Software}
    \label{fig:my_product}
\end{figure}
\section{Testphase}
\section{Einführung}
\section{Schlusswort}


\newpage

\setcounter{tocdepth}{1}
\listoffigures

\newpage
\printglossary[type=\acronymtype]
\newpage
\printglossary


\end{document}
